
\documentclass[10pt]{article}
\usepackage[margin=0.72in]{geometry}
\usepackage{amsmath,amssymb,booktabs,array,longtable,graphicx,xcolor,microtype}
\usepackage[hidelinks]{hyperref}
\usepackage{fancyhdr}
\usepackage{enumitem}
\graphicspath{{../figures/fullobs/}}
\definecolor{mottblue}{HTML}{335F8D}
\definecolor{chargeorange}{HTML}{C56F2D}
\definecolor{softgray}{HTML}{F4F5F7}
\pagestyle{fancy}
\fancyhf{}
\lhead{U=40 full-observable scan}
\rhead{Technical report}
\cfoot{\thepage}
\setlength{\headheight}{13pt}
\setlength{\parindent}{0pt}
\setlength{\parskip}{5pt}
\setlist[itemize]{leftmargin=1.3em,itemsep=2pt,topsep=2pt}
\newcommand{\vk}{\mathbf{k}}

\begin{document}

\begin{center}
{\LARGE\bfseries Static observables expose the Mott--charge reorganization at $U=40$\par}
\vspace{0.35em}
{\large All-momentum analysis of the $V$-extended MMHK calculation\par}
\vspace{0.4em}
{\small Data snapshot: 17 July 2026, 12:32 EDT}
\end{center}

\vspace{0.5em}
\colorbox{softgray}{\parbox{0.965\linewidth}{
\textbf{Technical summary.}
The analysis uses 7736 validated outputs out of 7749 requested
fixed-$Q=(0,0)$ calculations (99.832\% complete) at $U=40$, $t'=0$.
Every real-space charge and spin diagnostic changes branch between $V=10$ and $V=11$:
the on-site double occupation rises from 0.005155 to 0.497100, the checkerboard
order proxy rises from 0.001297 to 0.989141, and
$S^{zz}(\pi,\pi)$ collapses from 1.455514 to 0.002888.
The momentum distribution remains almost perfectly linear in the bare dispersion
($R^2\ge 0.999713$), but its slope changes from the Mott denominator $U-V$ to the
charge denominator $7V-U$.  A local Hellmann--Feynman branch interpolation gives
$V_c\simeq 10.0153$, or $V_c/U\simeq 0.250384$, while the directly sampled statement
is the robust bracket $10<V_c<11$.  The static observables therefore provide an
independent, sharply resolved counterpart to the Green-function-zero reorganization.}}

\section{A collective branch change occurs between $V=10$ and $V=11$}

The scan separates into two internally uniform regimes.  On the Mott side through
$V=10$, $D_i$ remains small, $\langle n_i n_j\rangle_{\rm NN}$ remains close to one,
and antiferromagnetic correlations remain strong.  At $V=11$, the same calculation
has instead entered an almost classical checkerboard state: $D_i\approx 1/2$,
$\langle n_i n_j\rangle_{\rm NN}\approx0$, and $m_{\rm CDW}^2\approx1$.
The numerical changes from $V=10$ to 11 are
\begin{align*}
D_i &: 0.005155 \rightarrow 0.497100, &
\langle n_i n_j\rangle_{\rm NN} &: 0.997420 \rightarrow 0.010137,\\
C_1 &: -0.345285 \rightarrow -0.001083, &
S^{zz}(\pi,\pi) &: 1.455514 \rightarrow 0.002888,\\
m_{\rm CDW}^2 &: 0.001297 \rightarrow 0.989141, &
K/N &: -0.312914 \rightarrow -0.214738.
\end{align*}
The momentum-to-momentum standard deviations of the local charge observables are
only $10^{-5}$ or smaller at the two boundary points.  Thus the switch is not caused
by a small subset of twist points; all sampled blocks select the same side at each
integer $V$.

\begin{figure}[htbp]
\centering
\includegraphics[width=\linewidth]{fullobs_static_scan_U40.pdf}
\caption{Twist-averaged static observables versus $V/U$.  Blue and orange backgrounds
mark the Mott- and charge-dominated sides of the atomic scale $V/U=1/4$.
The abrupt and simultaneous charge, spin, kinetic, and momentum changes occur between
the last Mott point ($V=10$) and the first charge point ($V=11$).}
\end{figure}

\section{$n(\vk)$ measures virtual hopping, not a generic fluctuation amplitude}

The full momentum maps show a simple but informative result.  For every $V$,
\begin{equation}
n(\vk)=1+a\,\epsilon_{\vk}+\delta n_{\vk},\qquad
\epsilon_{\vk}=-2(\cos k_x+\cos k_y),
\end{equation}
with $R^2\ge 0.999713$ and a very small nonlinear remainder.  Consequently the
momentum contrast $\sigma_{\vk}[n]$ and kinetic energy are not independent:
for the present normalization, $\sigma_{\vk}[n]=2|a|$ and $K/N=4a$ to the numerical
precision of the complete grids.

The momentum differentiation grows by 31.7\% from
$\sigma_{\vk}[n]=0.118794$ at $V=0$ to 0.156479 at $V=10$.
It then falls by -31.4\% to 0.107371 at $V=11$ and decreases further to
0.040132 at $V=20$.  Thus $n(\vk)$ is a clean measure of the virtual-hopping
amplitude and band-memory retained by the insulating state.  It is largest immediately
on the Mott side because the virtual charge denominator $U-V$ is smallest there,
then becomes progressively flatter in the charge state as $7V-U$ grows.

\begin{figure}[p]
\centering
\includegraphics[width=0.94\linewidth]{fullobs_nk_maps_U40.pdf}
\caption{Reconstructed $n(\vk)$ maps on $0\le k_x,k_y\le\pi$.  The Mott-side
dispersion strengthens toward $V/U=1/4$, drops immediately on the charge side,
and then flattens continuously.  Gray cells at $V/U=0.4$ and 0.5 denote the 13
still-unavailable calculations; no interpolation was used.}
\end{figure}

\begin{figure}[p]
\centering
\includegraphics[width=0.94\linewidth]{fullobs_nk_vs_epsilon_U40.pdf}
\caption{The measured momentum occupation versus the bare dispersion.  Orange lines
are weighted fits constrained by particle--hole symmetry to pass through
$n=1$ at $\epsilon_{\vk}=0$.  The near-unity $R^2$ values demonstrate that a
single interaction-dependent virtual-hopping scale describes each full map.}
\end{figure}

The separately measured momentum-orbital double occupation sharpens the interpretation.
Its average drops from $D_{\vk}=0.286465$ at $V=10$ to 0.252962 at $V=11$,
while the spin covariance
$D_{\vk}-[n(\vk)/2]^2$ collapses from 3.034e-02 to 8.032e-05.
However, the intrinsic number variance
$\mathrm{Var}(n_{\vk})=n(\vk)+2D_{\vk}-n(\vk)^2$ changes from 0.548445 to
0.494396 and approaches 0.499220 in the charge limit, rather than vanishing.
A flat $n(\vk)$ therefore means loss of momentum differentiation; it does not by
itself mean that the symmetry-preserving many-body state has zero occupation-number
variance.

\begin{figure}[p]
\centering
\includegraphics[width=0.94\linewidth]{fullobs_Dk_maps_U40.pdf}
\caption{Momentum-orbital double occupation
$D_{\vk}=\langle n_{\vk\uparrow}n_{\vk\downarrow}\rangle$.  The charge side tends
toward $D_{\vk}=1/4$, causing the up--down momentum covariance to vanish even while
the real-space double occupation tends to $1/2$.}
\end{figure}

\section{Two virtual charge denominators explain the measured trends}

The same nearest-neighbor hop has different excitation energies in the two atomic
backgrounds.  In a singly occupied Mott background it creates a neighboring
doublon--holon pair with
\begin{equation}
\Delta_M=U-V.
\end{equation}
In the checkerboard doublon--holon background it creates two adjacent singly occupied
sites.  The changed site touches seven interaction bonds, giving
\begin{equation}
\Delta_C=7V-U.
\end{equation}
Second-order perturbation theory at $t=1$ then gives the following parameter-free
forms once the measured nearest-neighbor spin correlation is inserted on the Mott side:
\begin{align}
\frac{E_M}{N} &\simeq 2V-\frac{8(1/4-C_1)}{\Delta_M}, &
\frac{K_M}{N} &\simeq -\frac{16(1/4-C_1)}{\Delta_M},\\
n_M(\vk) &\simeq 1-\frac{4(1/4-C_1)}{\Delta_M}\epsilon_{\vk}, &
D_M &\simeq \frac{8(1/4-C_1)}{\Delta_M^2},\\
\langle n_i n_j\rangle_M &\simeq 1-\frac{4(1/4-C_1)}{\Delta_M^2},\nonumber
\end{align}
and
\begin{align}
\frac{E_C}{N} &\simeq \frac{U}{2}-\frac{4}{\Delta_C}, &
\frac{K_C}{N} &\simeq -\frac{8}{\Delta_C},\\
n_C(\vk) &\simeq 1-\frac{2}{\Delta_C}\epsilon_{\vk}, &
D_C &\simeq \frac12-\frac{4}{\Delta_C^2},\\
\langle n_i n_j\rangle_C &\simeq \frac{14}{\Delta_C^2}.\nonumber
\end{align}
Across all applicable points, the $n(\vk)$ slope differs from these formulas by at
most 1.44\% on the Mott side and 0.68\% on the
charge side.  The largest Mott-side error in $D_i$ is only 2.58\%.
This agreement explains both why momentum differentiation grows toward the boundary
from the Mott side and why it decreases monotonically after the branch switch.

\begin{figure}[htbp]
\centering
\includegraphics[width=\linewidth]{fullobs_strong_coupling_validation_U40.pdf}
\caption{Measured energies, $n(\vk)$ slopes, kinetic energies, and local charge
probabilities compared with the two strong-coupling rules.  Each asymptotic rule is
shown only on its corresponding side.  The agreement is quantitative even for the
points adjacent to the branch change.}
\end{figure}

The exact Hellmann--Feynman relation for these outputs is
\begin{equation}
\frac{\partial(E_0/N)}{\partial V}=2\langle n_i n_j\rangle_{\rm NN}.
\end{equation}
Approximating each branch by its local tangent at $V=10$ and 11 and integrating over
the unit interval locates their crossing at $V_c\simeq 10.0153$,
$V_c/U\simeq 0.250384$.  The estimate shows that the small kinetic correction
barely shifts the atomic value.  It is nevertheless an interpolation; the direct
calculation establishes only the interval $10<V_c<11$.

\begin{figure}[htbp]
\centering
\includegraphics[width=\linewidth]{fullobs_transition_diagnostics_U40.pdf}
\caption{Left: the energy secants compared with the Hellmann--Feynman derivative.
The slope discontinuity identifies a branch crossing.  Right: order diagnostics on
a logarithmic scale, displaying the simultaneous loss of spin response and onset of
checkerboard charge order.}
\end{figure}

\section{Scope, definitions, and aggregation}

The calculation fixes $N_\uparrow=N_\downarrow=8$, $U=40$, $t=1$, $t'=0$, and
$Q=(0,0)$ at a $4\times4$ momentum-mixing resolution.  For each interaction pair,
the Peierls/twist momentum is sampled at
$f_{x,y}=k_{x,y}\pi/40+10^{-5}$ with $0\le k_x\le k_y\le40$.
The stored wedge is reflected under $x\leftrightarrow y$ to reconstruct the square;
off-diagonal rows have symmetry weight two.  Reported momentum averages additionally
use half weight at the $0$ and $\pi$ endpoints, i.e. two-dimensional trapezoidal
quadrature.  Missing rows are not imputed; each mean is normalized by the observed
quadrature weight.

The principal observables are:
\begin{itemize}
\item $D_i=\langle n_{i\uparrow}n_{i\downarrow}\rangle$ and
$D_{\vk}=\langle n_{\vk\uparrow}n_{\vk\downarrow}\rangle$;
\item $S^{zz}(\pi,\pi)=N^{-1}\sum_{ij}(-1)^{i-j}\langle S_i^zS_j^z\rangle$ and
$S_c(\pi,\pi)=N^{-1}\sum_{ij}(-1)^{i-j}\langle(n_i-1)(n_j-1)\rangle$;
\item $m_{\rm CDW}^2=S_c(\pi,\pi)/N$, the nearest-neighbor spin correlation $C_1$,
the nearest-neighbor density product, and kinetic energy per site;
\item $n(\vk)$ at $q=(0,0)$ of each $f$-twisted block, identified with physical
momentum $\vk=f$ in this construction.
\end{itemize}
The output field called \texttt{Gap} is identically zero because only one eigenpair
was requested.  It is metadata, not a charge, spin, or sector gap, and is excluded
from all physical conclusions.  The next-nearest-neighbor kinetic contribution is
also identically zero because $t'=0$.

\section{Numerical validation and incomplete-row sensitivity}

Every downloaded file is marked \texttt{COMPLETE}, matches its SHA-256 manifest,
and has a unique $(U,V,k_x,k_y)$ key.  Several identities provide independent checks:
\begin{align*}
E_0/N &= K/N+UD_i+2V\langle n_i n_j\rangle_{\rm NN},\\
\langle\delta n_i\delta n_j\rangle_{\rm NN}
&=\langle n_i n_j\rangle_{\rm NN}-1,\\
K&=K_{\rm NN}+K_{\rm NNN},\qquad
m_{\rm CDW}^2=S_c(\pi,\pi)/N.
\end{align*}
The largest energy-identity residual is 2.17e-10; the other three identities
hold to machine precision.  Particle--hole symmetry gives
$n(\vk)+n(\pi-\vk)=2$ with maximum residual below $7\times10^{-6}$.

All sets through $V=12$ are complete.  The $V=16$ and $V=20$ sets have 859/861 and
850/861 wedge rows, respectively.  Even assigning every missing $V=20$ value to the
observed minimum or maximum changes the mean $n(\vk)$ by at most 1.05e-03 and
$m_{\rm CDW}^2$ by at most 7.74e-09.  The missing rows therefore do not affect the
regime classification or strong-coupling conclusions.

\begin{table}[htbp]
\centering
\caption{Downloaded data coverage.  Quadrature coverage includes endpoint weights.}
\begin{tabular}{rrrr}
\toprule
$V$ & wedge rows & row coverage & quadrature coverage\\
\midrule
0 & 861 & 100.000\% & 100.000\% \\
4 & 861 & 100.000\% & 100.000\% \\
8 & 861 & 100.000\% & 100.000\% \\
9 & 861 & 100.000\% & 100.000\% \\
10 & 861 & 100.000\% & 100.000\% \\
11 & 861 & 100.000\% & 100.000\% \\
12 & 861 & 100.000\% & 100.000\% \\
16 & 859 & 99.768\% & 99.750\% \\
20 & 850 & 98.722\% & 98.688\% \\
\bottomrule
\end{tabular}
\end{table}

\section{Limitations and recommended next calculations}

\begin{itemize}
\item The direct interaction resolution brackets the switch only between $V=10$ and
11.  The value $V_c\simeq 10.0153$ is a local branch interpolation, not a separately
computed transition point.
\item The scan fixes $U=40$ and $Q=(0,0)$.  It establishes the strong-coupling
reorganization within the present momentum-mixing construction but does not by
itself provide thermodynamic scaling in mixing resolution or an unrestricted
momentum-sector comparison.
\item $n(\vk)$ contrast is an excellent virtual-hopping diagnostic, but it should be
called momentum differentiation rather than an undifferentiated measure of quantum
fluctuation.  The separately derived number variance remains finite in the charge
state.
\end{itemize}

The most informative next run is a narrow fixed-$U=40$ grid at
$V=10.02,10.05,10.10,$ and $10.25$, rather than another broad scan.  A second priority
is repeating a minimal set on both sides at a different large $U$ to test whether the
width and location collapse in $V/U$.  For the manuscript, the immediate next analysis
is to compare the measured $n(\vk)$ slope with the Green-function-zero scale at the
same $(U,V)$ points; the two observables are controlled by the same virtual charge
denominators and may provide a unified interpretation.

\section*{Further questions}
\begin{itemize}
\item Does the zero-dispersion scale inherit the same $U-V$ and $7V-U$ denominators
seen directly in $n(\vk)$ and $K$?
\item Does the branch switch remain as sharp when the momentum-mixing resolution or
total-momentum sector is varied?
\item Can the small nonlinear remainder $\delta n_{\vk}$ isolate higher-order spin
or charge correlations beyond the leading strong-coupling rules?
\end{itemize}

\section*{Reproducibility artifacts}
The frozen raw snapshot is
\path{/Users/cosdis/Desktop/projects/HK_orbital_cluster/benchmark/data/fullobs/U40_all_twist_20260717_123231}.  The parser, validation, and figure workflow is
\path{scripts/analyze_fullobs_u40_scan.py}; the report generator is
\path{scripts/generate_fullobs_u40_report.py}.  Machine-readable row data, long-form
observable statistics, wide summaries, the observable dictionary, and validation
JSON are stored under \path{results/} with tag \texttt{20260717\_123231}.

\appendix
\section{Complete numerical summaries}

\begin{table}[!htbp]
\centering\scriptsize
\caption{Energy, kinetic, and nearest-neighbor observables.}
\begin{tabular}{rrrrrrr}
\toprule
$V/U$ & $E_0/N$ & $K/N$ & $D_i$ & $\langle n_i n_j\rangle$ &
$\langle\delta n_i\delta n_j\rangle$ & $C_1$\\
\midrule
0.000 & -0.11948 & -0.237576 & 0.002952 & 0.998534 & -0.001466 & -0.347679 \\
0.100 & 7.86753 & -0.262878 & 0.003620 & 0.998200 & -0.001800 & -0.346953 \\
0.200 & 15.85136 & -0.294202 & 0.004547 & 0.997731 & -0.002269 & -0.345946 \\
0.225 & 17.84668 & -0.303258 & 0.004836 & 0.997584 & -0.002416 & -0.345631 \\
0.250 & 19.84169 & -0.312914 & 0.005155 & 0.997420 & -0.002580 & -0.345285 \\
0.275 & 19.89226 & -0.214738 & 0.497100 & 0.010137 & -0.989863 & -0.001083 \\
0.300 & 19.90933 & -0.180882 & 0.497947 & 0.007181 & -0.992819 & -0.000768 \\
0.400 & 19.94450 & -0.110879 & 0.499231 & 0.002692 & -0.997308 & -0.000288 \\
0.500 & 19.96002 & -0.079912 & 0.499601 & 0.001398 & -0.998602 & -0.000150 \\
\bottomrule
\end{tabular}
\end{table}

\begin{table}[!htbp]
\centering\scriptsize
\caption{Spin and charge-order observables.}
\begin{tabular}{rrrrrr}
\toprule
$V/U$ & $S^{zz}(\pi,\pi)$ & $S_c(\pi,\pi)$ & $m_{\rm CDW}^2$ & $D_i$ &
$\langle n_i n_j\rangle$\\
\midrule
0.000 & 1.464837 & 0.011728 & 0.000733 & 0.002952 & 0.998534 \\
0.100 & 1.462082 & 0.014423 & 0.000901 & 0.003620 & 0.998200 \\
0.200 & 1.458158 & 0.018217 & 0.001139 & 0.004547 & 0.997731 \\
0.225 & 1.456908 & 0.019419 & 0.001214 & 0.004836 & 0.997584 \\
0.250 & 1.455514 & 0.020758 & 0.001297 & 0.005155 & 0.997420 \\
0.275 & 0.002888 & 15.826261 & 0.989141 & 0.497100 & 0.010137 \\
0.300 & 0.002047 & 15.876917 & 0.992307 & 0.497947 & 0.007181 \\
0.400 & 0.000768 & 15.953847 & 0.997115 & 0.499231 & 0.002692 \\
0.500 & 0.000399 & 15.976040 & 0.998502 & 0.499601 & 0.001398 \\
\bottomrule
\end{tabular}
\end{table}

\begin{table}[!htbp]
\centering\scriptsize
\caption{Momentum observables and fitted band-memory scale.  The standard deviation
is over momentum, not a sampling error.}
\begin{tabular}{rrrrrrrr}
\toprule
$V/U$ & $\langle n\rangle$ & $\sigma_{\vk}[n]$ & $\langle D_{\vk}\rangle$ &
$\langle\mathrm{Var}(n_{\vk})\rangle$ & $\langle\mathrm{Cov}_{\uparrow\downarrow}\rangle$ & $a$ & $R^2$\\
\midrule
0.000 & 0.9999985 & 0.118794 & 0.284419 & 0.554727 & 3.09e-02 & -0.059394 & 0.999889 \\
0.100 & 0.9999983 & 0.131449 & 0.285053 & 0.552830 & 3.07e-02 & -0.065719 & 0.999848 \\
0.200 & 0.9999981 & 0.147118 & 0.285913 & 0.550184 & 3.05e-02 & -0.073551 & 0.999775 \\
0.225 & 0.9999981 & 0.151648 & 0.286176 & 0.549357 & 3.04e-02 & -0.075815 & 0.999747 \\
0.250 & 0.9999980 & 0.156479 & 0.286465 & 0.548445 & 3.03e-02 & -0.078228 & 0.999713 \\
0.275 & 0.9999986 & 0.107371 & 0.252962 & 0.494396 & 8.03e-05 & -0.053684 & 0.999955 \\
0.300 & 0.9999989 & 0.090442 & 0.252100 & 0.496021 & 5.56e-05 & -0.045221 & 0.999978 \\
0.400 & 1.0000968 & 0.055460 & 0.250840 & 0.498507 & 2.22e-05 & -0.027720 & 0.999997 \\
0.500 & 1.0001541 & 0.040132 & 0.250492 & 0.499220 & 1.25e-05 & -0.019978 & 0.999999 \\
\bottomrule
\end{tabular}
\end{table}

\begin{figure}[p]
\centering
\includegraphics[width=0.94\linewidth]{fullobs_varnk_maps_U40.pdf}
\caption{Derived intrinsic momentum-orbital number variance.  Its finite charge-side
limit distinguishes occupation-number variance from the flattening of the mean
momentum distribution.  Gray cells are unavailable raw rows.}
\end{figure}

\end{document}
